% Краткий справочник формул
% Лаборатория математики ФН1
% Компиляция: pdflatex spravochnik-formul.tex

\documentclass[11pt, a4paper]{article}

\usepackage[T2A]{fontenc}
\usepackage[utf8]{inputenc}
\usepackage[russian]{babel}
\usepackage[left=18mm, right=18mm, top=18mm, bottom=18mm]{geometry}
\usepackage{amsmath, amssymb}
\usepackage{multicol}
\usepackage{array}

\pagestyle{plain}
\setlength{\parindent}{0pt}
\setlength{\columnsep}{8mm}

\newcommand{\block}[1]{\vspace{0.6em}\textbf{#1}\vspace{0.25em}\hrule\vspace{0.5em}}

\begin{document}

\begin{center}
  {\Large\bfseries Краткий справочник формул}

  \vspace{0.3em}
  {\small Лаборатория математики ФН1 · для подготовки к контрольным и экзаменам}
\end{center}

\vspace{0.8em}

\begin{multicols}{2}

\block{Производные}

\begin{tabular}{@{}ll@{}}
  $(x^{\alpha})' = \alpha x^{\alpha-1}$ & $(e^x)' = e^x$ \\[2pt]
  $(a^x)' = a^x\ln a$ & $(\ln x)' = \dfrac{1}{x}$ \\[4pt]
  $(\sin x)' = \cos x$ & $(\cos x)' = -\sin x$ \\[4pt]
  $(\operatorname{tg} x)' = \dfrac{1}{\cos^2 x}$ & $(\operatorname{ctg} x)' = -\dfrac{1}{\sin^2 x}$ \\[6pt]
  $(\arcsin x)' = \dfrac{1}{\sqrt{1-x^2}}$ & $(\operatorname{arctg} x)' = \dfrac{1}{1+x^2}$
\end{tabular}

\vspace{0.4em}
Правила: $(uv)' = u'v + uv'$, \quad
$\left(\dfrac{u}{v}\right)' = \dfrac{u'v - uv'}{v^2}$, \quad
$\big(f(g)\big)' = f'(g)\,g'$.

\block{Неопределённые интегралы}

\begin{tabular}{@{}ll@{}}
  $\displaystyle\int x^{\alpha}dx = \frac{x^{\alpha+1}}{\alpha+1}$ & $\alpha \ne -1$ \\[8pt]
  $\displaystyle\int \frac{dx}{x} = \ln|x|$ & \\[8pt]
  $\displaystyle\int e^x dx = e^x$ & \\[8pt]
  $\displaystyle\int \frac{dx}{x^2 + a^2} = \frac{1}{a}\operatorname{arctg}\frac{x}{a}$ & \\[8pt]
  $\displaystyle\int \frac{dx}{\sqrt{a^2 - x^2}} = \arcsin\frac{x}{a}$ &
\end{tabular}

\vspace{0.4em}
По частям: $\displaystyle\int u\,dv = uv - \int v\,du$.

\block{Замечательные пределы}

$\displaystyle\lim_{x\to 0}\frac{\sin x}{x} = 1$, \qquad
$\displaystyle\lim_{x\to\infty}\left(1 + \frac{1}{x}\right)^{x} = e$.

\block{Разложения при $x \to 0$}

$e^x = 1 + x + \dfrac{x^2}{2!} + \dots$

$\sin x = x - \dfrac{x^3}{3!} + \dfrac{x^5}{5!} - \dots$

$\cos x = 1 - \dfrac{x^2}{2!} + \dfrac{x^4}{4!} - \dots$

$\ln(1+x) = x - \dfrac{x^2}{2} + \dfrac{x^3}{3} - \dots$

$(1+x)^{\alpha} = 1 + \alpha x + \dfrac{\alpha(\alpha-1)}{2!}x^2 + \dots$

\columnbreak

\block{Матрицы и определители}

$\det(AB) = \det A \cdot \det B$, \quad $\det A^{\mathsf T} = \det A$

$A^{-1} = \dfrac{1}{\det A}\,\widetilde{A}^{\mathsf T}$, \quad $\det A \ne 0$

$(AB)^{-1} = B^{-1}A^{-1}$, \quad $(AB)^{\mathsf T} = B^{\mathsf T}A^{\mathsf T}$

Теорема Кронекера~--- Капелли: система совместна тогда и только тогда,
когда $\operatorname{rk} A = \operatorname{rk}(A\,|\,b)$.

\block{Собственные значения}

$\det(A - \lambda E) = 0$, \quad
$\sum \lambda_i = \operatorname{tr} A$, \quad
$\prod \lambda_i = \det A$.

\block{Теория вероятностей}

$P(A \cup B) = P(A) + P(B) - P(AB)$

$P(A\,|\,B) = \dfrac{P(AB)}{P(B)}$

Полная вероятность: $P(A) = \sum_i P(H_i)P(A\,|\,H_i)$

Байес: $P(H_k\,|\,A) = \dfrac{P(H_k)P(A\,|\,H_k)}{P(A)}$

Бернулли: $P_n(k) = C_n^k p^k q^{n-k}$

$M[X] = \sum x_i p_i$, \quad
$D[X] = M[X^2] - \big(M[X]\big)^2$

\block{Численные методы}

Метод Ньютона: $x_{n+1} = x_n - \dfrac{f(x_n)}{f'(x_n)}$

Формула трапеций: $\displaystyle\int_a^b f\,dx \approx h\left(\frac{f_0 + f_n}{2} + \sum_{i=1}^{n-1} f_i\right)$

Формула Симпсона: погрешность $O(h^4)$

Наблюдаемый порядок: $p = \log_2 \dfrac{\varepsilon_h}{\varepsilon_{h/2}}$

\end{multicols}

\vfill
{\footnotesize Справочник не заменяет конспект: формулы даны без условий
применимости. Проверяйте их в основной литературе раздела.}

\end{document}
